\section{预备记号与对象层级}

\subsection{微态读出与窗口词}
固定分辨率 $m\in\NN_{\ge 1}$，定义微态空间
$$
\Omega_m := \{0,1\}^m.
$$
微态词可视为某一基础系统经窗口读出的长度 $m$ 二元序列。本文不将底层系统作为原语；POM 的最小接口仅要求存在“读出可被地址触发并生成有限串”的过程，并以此作为对象生成的输入层。

\subsection{地址、柱事件与空值语义}
令 $\Addr$ 为可数地址集合（例如有限二元串、带参数的投影词等）。对地址 $a\in\Addr$，定义其触发的读出过程 $\Read(a)$。若地址未给定，则读出过程不定义，记为 $\NULL$，并约定 $\NULL$ 不属于数值域且 $\NULL\neq 0$。该“地址先于值”的语义是 POM 的存在性基础。

\subsection{稳定语法与稳定对象}
定义稳定类型集合 $X_m\subseteq\Omega_m$ 为满足局部语法约束的子集。本文重点讨论黄金均值语法（禁止相邻 $11$）：
$$
X_m := \Bigl\{x\in\{0,1\}^m:\ \text{不存在相邻位置同时为 }1\Bigr\}.
$$
其逆极限对象 $X_\infty$ 可视为无限合法序列空间（并在有限支撑语义下等价于 Zeckendorf 合法数字空间）。POM 将“外表可见对象”定义为落在稳定层（稳定子代数）中的可观测量；该可见性在后续实现上对应于被条件期望族 $E_m$ 不变的观测层。

